% Entry points into a vehicle

\begin{frame}{How Do You Break Into a Car?}
  \begin{itemize}
    \item \textbf{Touch it}: OBD-II port, USB, the EV charging port, a physical CAN tap.
    \item \textbf{Stand near it}: Bluetooth, Wi-Fi, the key fob (relay/replay), tyre-pressure sensors.
    \item \textbf{From anywhere}: the cellular telematics link and the cloud APIs that run the whole fleet.
    \item \textbf{Never touch it}: the supply chain, open-source libraries, supplier firmware, dealer flashing tools.
    \item \textbf{The update pipe itself}: OTA servers, signing keys, the CI/CD that builds the car's software.
  \end{itemize}
  \vspace{1mm}
  \begin{alertblock}{The uncomfortable truth}
    The weakest link is often not the CAN bus or the ECU firmware, it's the backend that controls millions of cars. In 2022, web-API bugs exposed 16+ carmakers at once. Automotive security increasingly includes product security: cloud APIs, mobile apps, update systems and backend services.
  \end{alertblock}
\end{frame}
